\section{Términos clave: índice de palabras clave de este episodio}

\begin{longtable}{p{0.24\textwidth}p{0.32\textwidth}p{0.35\textwidth}}
\toprule
Término & Explicación en una frase & Papel en este episodio \\
\midrule
World Model & Modelo que predice estados futuros a partir del estado y la acción & Tesis técnica central de la entrevista. \\
Predictive Brain & Cerebro predictivo orientado al mundo físico & Capa de inteligencia de nivel superior que AMI Labs quiere construir. \\
State & Representación comprimida del entorno, relevante para la tarea & Vincula el aprendizaje de representaciones con el planning. \\
Action & Acción o intervención que el agente ejerce sobre el entorno & Hace que la predicción intervenga en el control y la toma de decisiones. \\
Representation Learning & Aprendizaje de representaciones abstractas reutilizables & Puente que lleva de la vision al world model. \\
Vision as Perspective & Entender la visión como la perspectiva fundamental de la inteligencia, no como un conjunto de tareas & Explica el sentido de la CV en la era de los LLM. \\
LLM-pilled & Que la narrativa de los LLM organiza en exceso los recursos y la imaginación & Crítica a la cadena de valor de Silicon Valley. \\
VLA & Modelo Vision-Language-Action & Forma importante, pero insuficiente, de la línea de la robótica. \\
World Model Pre-training & Preentrenamiento fundacional orientado a video, sensores y el mundo físico & Problema abierto central de la segunda mitad del preentrenamiento. \\
Research Taste & Capacidad de juzgar problemas, evidencia, momento oportuno y estética & Explica las decisiones de personas y organizaciones. \\
JEPA & Joint Embedding Predictive Architecture & Idea representativa de la línea de LeCun: «predecir en un espacio de representaciones abstractas». \\
OpenAI a la inversa & Construir modelos junto con el mundo y una red de socios, en lugar de solo descargar internet & Metáfora organizacional y comercial de AMI. \\
\bottomrule
\end{longtable}

\subsection{Resumen de la sección}

Todas estas palabras clave apuntan a la misma pregunta: si en la siguiente etapa la IA podrá pasar de «saber procesar el mundo que los humanos escribieron» a «poder predecir, actuar y aprender en el mundo real». Esto exige actualizar al mismo tiempo los modelos, los datos, la organización y los límites filosóficos.

